\documentclass[12pt]{article}

% O babel vem antes do modelo SBC: com o hyperref carregado, o bloco final do
% sbc-template.sty executa \makeatother no fim do preâmbulo.
\usepackage[utf8]{inputenc}
\usepackage[T1]{fontenc}
\usepackage[brazil]{babel}
\usepackage{sbc-template}
\usepackage{graphicx,url}
\usepackage{booktabs}
\usepackage{tabularx}
\usepackage{array}
\usepackage{multirow}
\usepackage{xcolor}
\usepackage{tikz}
\usetikzlibrary{positioning,arrows.meta,fit,calc}
\usepackage[hidelinks]{hyperref}

\sloppy
\newcolumntype{L}{>{\raggedright\arraybackslash}X}
\newcommand{\tabfonte}[1]{\par\vspace{2pt}{\footnotesize Fonte: #1}}
\renewcommand{\arraystretch}{1.15}

\title{Avaliação da Qualidade de Produto do Moodle 5.3 LTS:\\
Adequação Funcional, Compatibilidade e Flexibilidade\\
segundo a ISO/IEC 25010:2023 e o GQM}

\author{Bernardo Broetto Brun, Carlos Eduardo Deusdara Motta,\\
Maykon Junio dos Santos Soares, Pedro Goncalves Rocha}

\address{Faculdade de Ciências e Tecnologias em Engenharia (FCTE)\\
Universidade de Brasília (UnB) -- Gama, DF -- Brasil}

\begin{document}

\maketitle

\begin{center}\footnotesize
Repositório: \url{https://github.com/FCTE-Qualidade-de-Software-1/2026-2_T01_MARLYN_WESCOFF}\\
GitPage: \url{https://fcte-qualidade-de-software-1.github.io/2026-2_T01_MARLYN_WESCOFF/}
\end{center}

\begin{abstract}
This paper reports the first two phases of a product quality evaluation of Moodle 5.3 LTS, an open source learning management system, following the ISO/IEC 25040 evaluation process and the ISO/IEC 25010:2023 product quality model. The evaluation supports a higher education institution deciding whether to upgrade to the new long-term support release. A weighted decision matrix with a feasibility filter and sensitivity analysis prioritized Functional Suitability, Compatibility and Flexibility. Using the Goal--Question--Metric approach, the three characteristics were refined into three measurement goals, eleven questions with hypotheses, and twelve metrics derived from ISO/IEC 25023, each with scoring levels and explicit judgment rules. Project execution is tracked with PSM/CID measures.
\end{abstract}

\begin{resumo}
Este artigo relata as duas primeiras fases da avaliação da qualidade de produto do Moodle 5.3 LTS, sistema de gestão da aprendizagem de código aberto, seguindo o processo de avaliação da ISO/IEC 25040 e o modelo de qualidade de produto da ISO/IEC 25010:2023. A avaliação apoia uma instituição de ensino superior na decisão de atualizar para a nova versão de suporte estendido. Uma matriz de decisão ponderada, com filtro de viabilidade e análise de sensibilidade, priorizou Adequação Funcional, Compatibilidade e Flexibilidade. Com o método \textit{Goal--Question--Metric}, as três características foram desdobradas em três objetivos de medição, onze questões com hipóteses e doze métricas derivadas da ISO/IEC 25023, cada uma com níveis de pontuação e regras de julgamento explícitas. A execução do projeto é acompanhada com medidas do PSM/CID.
\end{resumo}

{\small\tableofcontents}

% ======================================================================
\section{Introdução}
\label{sec:intro}

Ambientes virtuais de aprendizagem tornaram-se infraestrutura essencial do ensino superior. Quando um desses sistemas falha em uma tarefa acadêmica, perde notas ao trocar dados com a secretaria ou não abre no celular do estudante, o prejuízo atinge diretamente o processo de ensino. O \textbf{Moodle} é o sistema de gestão da aprendizagem (LMS) de código aberto mais adotado em universidades brasileiras, inclusive na UnB.

Em 05/10/2026 foi lançado o Moodle 5.3, nova versão de suporte estendido (LTS). A versão exige PHP 8.3 ou superior, PostgreSQL 17 ou superior e só pode ser atualizada a partir da 4.4~\cite{moodle53}; além disso, desde a 5.1 o código acessível pela web foi movido para a pasta \texttt{public/}, o que obriga a reconfigurar o servidor web na atualização~\cite{moodle51upgrade}. Uma instituição que mantém o Moodle precisa, portanto, decidir \emph{se} e \emph{como} migrar.

Este trabalho avalia a qualidade de produto do Moodle 5.3 LTS para apoiar essa decisão. A avaliação segue o processo da ISO/IEC 25040~\cite{iso25040}, adota o modelo de qualidade da ISO/IEC 25010:2023~\cite{iso25010} e especifica as medições com o método GQM~\cite{basili1994gqm}. O projeto foi conduzido pelo Grupo Marlyn Wescoff na disciplina Qualidade de Software~1 da FCTE/UnB~\cite{ramos2026trabalho}.

Esta primeira versão do artigo cobre as Fases~1 e~2 do processo. A Seção~\ref{sec:referencial} apresenta o referencial teórico; a Seção~\ref{sec:metodologia}, a metodologia de desenvolvimento do projeto; a Seção~\ref{sec:fase1}, o estabelecimento dos requisitos de avaliação (Fase~1); a Seção~\ref{sec:fase2}, a especificação da avaliação (Fase~2); e a Seção~\ref{sec:consideracoes}, as considerações parciais. O Apêndice~\ref{ap:contribuicao} traz a tabela de contribuição da equipe e o Apêndice~\ref{ap:ia}, a declaração de uso de IA.

% ======================================================================
\section{Referencial teórico}
\label{sec:referencial}

\subsection{Série SQuaRE e o modelo de qualidade da ISO/IEC 25010:2023}

A série ISO/IEC 25000, conhecida como SQuaRE (\textit{Systems and software Quality Requirements and Evaluation}), organiza em divisões as normas de gestão da qualidade (2500n), modelos de qualidade (2501n), medição (2502n), requisitos (2503n) e avaliação (2504n)~\cite{ramos2026aulas}. A ISO/IEC 25010:2023 define o modelo de qualidade de produto com nove características: adequação funcional, eficiência de desempenho, compatibilidade, capacidade de interação, confiabilidade, segurança, manutenibilidade, flexibilidade e segurança operacional~\cite{iso25010}. Em relação à edição de 2011, a revisão renomeou \emph{portabilidade} para \emph{flexibilidade}, incluiu a subcaracterística \emph{escalabilidade} e renomeou \emph{usabilidade} para \emph{capacidade de interação}.

As medidas de qualidade de produto são descritas na ISO/IEC 25023, que propõe, para cada subcaracterística, medidas na forma de razões entre contagens, como a cobertura de implementação funcional e a intercambialidade de formatos de dados~\cite{iso25023}.

\subsection{Processo de avaliação da ISO/IEC 25040}

A ISO/IEC 25040 organiza a avaliação em quatro fases~\cite{iso25040,ramos2026aulas}: (1)~\emph{estabelecer os requisitos de avaliação}, definindo propósito, tipo de produto e modelo de qualidade; (2)~\emph{especificar a avaliação}, selecionando métricas, níveis de pontuação e critérios de julgamento; (3)~\emph{projetar a avaliação}, produzindo o plano com método, recursos e cronograma; e (4)~\emph{executar a avaliação}, obtendo as medidas, comparando-as com os critérios e julgando os resultados. O detalhamento deve permitir que a avaliação seja repetível e reprodutível.

\subsection{Classificação do tipo de produto}

A Fase~1 exige identificar o tipo de produto. Pressman~\cite{pressman2002} classifica o software por domínio de aplicação (sistema, tempo real, comercial, científico, embarcado, aplicações web, entre outros). A IEEE~1062~\cite{ieee1062} distingue o software COTS (\textit{Commercial Off-The-Shelf}), igual para todos os clientes; o MOTS (\textit{Modified Off-The-Shelf}), um produto-padrão adaptado a clientes; e o software totalmente desenvolvido sob encomenda.

\subsection{Goal--Question--Metric}

O GQM, proposto por Basili, Caldiera e Rombach~\cite{basili1994gqm}, parte do princípio de que a medição deve ser orientada a objetivos. Ele tem três níveis: o \emph{conceitual} (objetivo de medição), o \emph{operacional} (questões que caracterizam o objetivo) e o \emph{quantitativo} (métricas que respondem às questões). O objetivo é descrito por um template com cinco elementos: \emph{analisar} um objeto, \emph{para o propósito de} entender, melhorar ou controlar, \emph{com respeito a} um foco de qualidade, \emph{do ponto de vista de} quem deseja a avaliação, \emph{no contexto de} um ambiente. As \emph{abstraction sheets} registram foco de qualidade, hipóteses de \textit{baseline}, fatores de variação e seu impacto, e apoiam a formulação de hipóteses para cada questão~\cite{ramos2026aulas}.

\subsection{PSM e o guia CID}

O PSM (\textit{Practical Software and Systems Measurement}) é uma abordagem de medição dirigida por necessidades de informação, com forte ligação com a gestão de projetos~\cite{mcgarry2002psm}. Seu modelo de informação liga entidades e atributos a medidas básicas, medidas derivadas e indicadores, interpretados por modelos de análise e critérios de decisão. O guia CID (\textit{Continuous Iterative Development Measurement Framework}) adapta o PSM a projetos iterativos e lista, em sua Tabela~5, conceitos mensuráveis e medidas para categorias como Calendário e Progresso, Recursos e Custos e Desempenho dos Processos~\cite{psmcid2021}.

% ======================================================================
\section{Metodologia de desenvolvimento do projeto}
\label{sec:metodologia}

\subsection{Processo e divisão em entregas}

O projeto segue as quatro fases da ISO/IEC 25040 (Figura~\ref{fig:processo}). A Entrega~1 (E1, 26/10/2026) cobre as Fases~1 e~2; a Entrega~2 cobre as Fases~3 e~4, incluindo a ação de melhoria obrigatória.

\begin{figure}[!ht]
\centering
\begin{tikzpicture}[
  fase/.style={draw, rounded corners, fill=blue!6, align=center, minimum width=3.0cm, minimum height=1.2cm, font=\footnotesize},
  seta/.style={-{Stealth}, thick}]
\node[fase] (f1) {\textbf{Fase 1}\\Estabelecer\\requisitos};
\node[fase, right=0.45cm of f1] (f2) {\textbf{Fase 2}\\Especificar\\a avaliação};
\node[fase, right=0.45cm of f2, fill=gray!8] (f3) {\textbf{Fase 3}\\Projetar\\a avaliação};
\node[fase, right=0.45cm of f3, fill=gray!8] (f4) {\textbf{Fase 4}\\Executar\\a avaliação};
\draw[seta] (f1) -- (f2); \draw[seta] (f2) -- (f3); \draw[seta] (f3) -- (f4);
\node[draw, dashed, rounded corners, fit=(f1)(f2), inner sep=5pt, label={[font=\footnotesize]below:Entrega 1}] {};
\node[draw, dashed, rounded corners, fit=(f3)(f4), inner sep=5pt, label={[font=\footnotesize]below:Entrega 2}] {};
\end{tikzpicture}
\caption{Processo de avaliação adotado e sua divisão entre as entregas. Fonte: elaboração própria a partir da ISO/IEC 25040.}
\label{fig:processo}
\end{figure}

\subsection{Organização da equipe e ferramentas}

A equipe tem quatro integrantes. O trabalho foi dividido por característica de qualidade, com um responsável e um revisor para cada uma; as seções transversais (propósito, partes interessadas, ODS, gestão, artigo e GitPage) foram distribuídas entre os integrantes. As decisões coletivas são registradas em atas no repositório.

O repositório Git da disciplina concentra documentos, dados e o código-fonte deste artigo. Cada tarefa é uma \emph{issue} com estimativa de horas, associada a um \emph{milestone} que representa a sprint; a integração na \texttt{main} é feita por \emph{pull request} revisado por outro integrante. A GitPage, gerada com MkDocs Material a partir da pasta \texttt{docs/}, é publicada automaticamente a cada \emph{push}, e o artigo é compilado por um \emph{workflow} do GitHub Actions. Os diagramas são versionados como texto (Mermaid e TikZ).

\subsection{Gestão com PSM/CID}

A especificação do trabalho exige medidas de gestão em três categorias de informação. A Tabela~\ref{tab:psm} resume as medidas escolhidas na Tabela~5 do CID, Parte~1~\cite{psmcid2021}, e seus critérios de decisão. Cada medida foi descrita no template de instância do modelo de informação da Aula~09~\cite{ramos2026aulas}; as instâncias completas, os retratos datados das medições e as decisões de gestão estão na GitPage. Os dados são extraídos das \emph{issues} pela API do GitHub e de um registro de horas por integrante (\texttt{dados/gestao/horas.csv}), por um script versionado no repositório.

\begin{table}[!ht]
\centering
\caption{Medidas do PSM/CID usadas na gestão do projeto.}
\label{tab:psm}
\footnotesize
\begin{tabularx}{\linewidth}{p{2.4cm}p{2.9cm}p{2.6cm}L}
\toprule
\textbf{Categoria} & \textbf{Medida} & \textbf{Função} & \textbf{Critério de decisão} \\
\midrule
Calendário e Progresso & Planejado vs. concluído & $C/P \times 100$ & Menos de 50\% na metade da sprint: redistribuir tarefas; diferença $>10\%$ no fim: registrar causa e replanejar \\
Recursos e Custos & Esforço estimado vs. realizado & $R/E$; custo $= R \times$ R\$~15,00 & $R/E > 1{,}3$: revisar estimativas e dividir tarefas \\
Desempenho dos Processos & Tempo de ciclo & mediana$(t_f - t_c)$ & Issue sem atualização há mais de 3 dias: tratar bloqueio \\
\bottomrule
\end{tabularx}
\tabfonte{elaboração própria a partir do PSM/CID~\cite{psmcid2021}.}
\end{table}

% ======================================================================
\section{Fase 1: estabelecimento dos requisitos de avaliação}
\label{sec:fase1}

\subsection{Propósito e uso pretendido}

O cenário é o de uma instituição pública de ensino superior que já usa o Moodle e analisa a atualização para a 5.3 LTS. O cenário é inspirado no uso do Moodle em universidades brasileiras, mas nenhuma instituição solicitou formalmente a avaliação. A avaliação deve produzir evidências sobre três riscos: \emph{acadêmico} (tarefas essenciais falharem ou produzirem resultados errados), \emph{de integração} (perda de dados nas trocas com sistemas vizinhos ou interferência em serviços do mesmo servidor) e \emph{operacional} (falha na instalação ou atualização, ou nos navegadores e dispositivos usados). A Tabela~\ref{tab:decisoes} relaciona esses riscos às decisões apoiadas.

\begin{table}[!ht]
\centering
\caption{Decisões apoiadas pela avaliação.}
\label{tab:decisoes}
\footnotesize
\begin{tabularx}{\linewidth}{p{0.6cm}Lp{2.6cm}p{2.4cm}}
\toprule
\textbf{ID} & \textbf{Decisão} & \textbf{Quem decide} & \textbf{Característica} \\
\midrule
D1 & Adotar a 5.3 LTS só com o núcleo ou prever plugins para tarefas essenciais & Coordenação de TI e gestão acadêmica & Adequação funcional \\
D2 & Manter as integrações atuais ou planejar adaptações e testes de regressão & Equipe de TI e integradores & Compatibilidade \\
D3 & Seguir o procedimento oficial de instalação/atualização ou preparar procedimento próprio; definir ambientes suportados & Administradores de TI & Flexibilidade \\
D4 & Escolher a ação de melhoria da Entrega~2 & Equipe avaliadora e docente & Todas \\
\bottomrule
\end{tabularx}
\tabfonte{elaboração própria.}
\end{table}

Cada objetivo de medição termina em um julgamento de três níveis --- \emph{atende}, \emph{atende com ressalvas} ou \emph{não atende} --- que corresponde, respectivamente, a prosseguir com a migração, prosseguir com ações de mitigação ou não prosseguir sem correção documentada (Seção~\ref{sec:julgamento}).

\subsection{Requisitante e partes interessadas}

A avaliação tem dois requisitantes: a \emph{docente da disciplina}, que define regras, prazos e critérios (inclusive a exclusão da usabilidade), e a \emph{Coordenação de TI} da instituição, papel de comprador e operador no cenário, que toma as decisões D1--D3 e é o ponto de vista dos objetivos GQM. A Tabela~\ref{tab:stakeholders} resume as demais partes interessadas. As necessidades marcadas com (H) foram derivadas da documentação do Moodle e da experiência dos integrantes como usuários, sem entrevistas, e serão tratadas como hipóteses.

\begin{table}[!ht]
\centering
\caption{Partes interessadas, critério de sucesso e influência na avaliação.}
\label{tab:stakeholders}
\footnotesize
\begin{tabularx}{\linewidth}{p{3.1cm}LL}
\toprule
\textbf{Parte interessada (papel)} & \textbf{Critério de sucesso} & \textbf{Influência na avaliação} \\
\midrule
Docente (requisitante acadêmica) & Métricas com fonte e procedimento definidos & Exclui usabilidade; exige GQM e PSM \\
Coordenação de TI (comprador/operador) & Julgamento por característica para decidir D1--D3 & Ponto de vista dos objetivos G1--G3 \\
Gestão acadêmica (H) & Todas as tarefas essenciais cobertas e corretas & Prioriza adequação funcional \\
Administradores (H) (operadores) & Instalar e atualizar sem erro bloqueante e sem perda de dados & Prioriza instalabilidade e coexistência \\
Integradores (H) (desenvolvedores) & Formatos e serviços sem perda de dados & Prioriza interoperabilidade \\
Docentes e tutores (H) (usuários) & Notas e prazos calculados conforme configurado & Define casos de correção funcional \\
Estudantes (H) (usuários) & Tarefas concluídas nos navegadores e no celular & Define ambientes de adaptabilidade \\
Comunidade Moodle (fornecedora) & Relatos reproduzíveis de defeitos & Fornece documentação de referência \\
\bottomrule
\end{tabularx}
\tabfonte{elaboração própria. (H): necessidade hipotética, ainda não validada.}
\end{table}

\subsection{Tipo de produto e descrição estruturada}

O Moodle é uma aplicação web cliente-servidor, um monolito modular em PHP extensível por plugins, distribuído sob GPLv3. Na classificação de Pressman, enquadra-se como aplicação web de sistema de informação; na IEEE~1062, é um produto \textbf{COTS}, que pode se tornar MOTS quando a instituição instala plugins e customizações~\cite{pressman2002,ieee1062}. Por isso, a avaliação usa apenas o núcleo, com o tema padrão Boost, para que os resultados não dependam de customizações locais.

Os módulos relevantes são cursos e matrículas, usuários (incluindo o cadastro em lote por CSV), tarefa (\texttt{mod\_assign}), questionário e banco de questões (\texttt{mod\_quiz}), fórum, recurso, livro de notas, calendário, backup e restauração, serviços web REST, instalação/atualização pela linha de comando e tarefas agendadas~\cite{hunt2012moodle}. As interfaces são a interface web, a API REST, os arquivos de importação e exportação (CSV, XLSX, GIFT, Moodle XML, iCal, \texttt{.mbz}), a linha de comando e o acesso ao banco de dados e ao diretório \texttt{moodledata}. A Figura~\ref{fig:contexto} mostra o contexto da instalação de referência.

\begin{figure}[!ht]
\centering
\begin{tikzpicture}[
  caixa/.style={draw, rounded corners, align=center, font=\scriptsize, minimum height=0.8cm, inner sep=3pt},
  ext/.style={caixa, fill=orange!8},
  int/.style={caixa, fill=blue!6},
  seta/.style={{Stealth}-{Stealth}, semithick}]
\node[int, minimum width=4.6cm, minimum height=2.3cm] (moodle) {};
\node[font=\scriptsize\bfseries, anchor=north] at (moodle.north) {Moodle 5.3 LTS (núcleo)};
\node[int, font=\tiny, anchor=south west] at ($(moodle.south west)+(0.12,0.12)$) (mods) {Tarefa, questionário,\\fórum, notas, calendário};
\node[int, font=\tiny, anchor=south east] at ($(moodle.south east)+(-0.12,0.12)$) (ws) {Serviços web REST\\backup, cron, CLI};
\node[ext, left=1.3cm of moodle, yshift=0.6cm] (us) {Estudantes e\\docentes (navegador)};
\node[ext, left=1.3cm of moodle, yshift=-0.7cm] (adm) {Administradores\\(CLI)};
\node[ext, right=1.3cm of moodle, yshift=0.75cm] (sis) {Sistema acadêmico\\(CSV, XLSX)};
\node[ext, right=1.3cm of moodle, yshift=-0.1cm] (cli) {Cliente REST};
\node[ext, right=1.3cm of moodle, yshift=-0.9cm] (cal) {Calendário (iCal)};
\node[int, below=0.75cm of moodle, xshift=-3.1cm] (php) {Apache 2.4 -- raiz \texttt{public/}\\PHP 8.3 / 8.4};
\node[int, below=0.75cm of moodle] (db) {PostgreSQL 17\\\texttt{moodledata}};
\node[int, below=0.75cm of moodle, xshift=3.1cm, fill=gray!10] (viz) {Aplicação vizinha\\(coexistência)};
\draw[seta] (us) -- (moodle.west |- us);
\draw[seta] (adm) -- (moodle.west |- adm);
\draw[seta] (sis) -- (moodle.east |- sis);
\draw[seta] (cli) -- (moodle.east |- cli);
\draw[seta] (cal) -- (moodle.east |- cal);
\draw[seta] (php) -- (php |- moodle.south);
\draw[seta] (db) -- (moodle.south);
\draw[semithick, dashed] (viz) -- (php);
\node[draw, dashed, rounded corners, fit=(php)(db)(viz), inner sep=4pt, label={[font=\tiny]below:Servidor de referência: Ubuntu Server 24.04 (máquina virtual)}] {};
\end{tikzpicture}
\caption{Contexto e arquitetura simplificada da instalação de referência. Fonte: elaboração própria a partir de~\cite{moodle53}.}
\label{fig:contexto}
\end{figure}

\textbf{Implicações para a avaliação.} Com a instalação local e dados fictícios, é possível medir cobertura e correção das tarefas essenciais, trocas de dados por arquivos e serviços web, coexistência com outra aplicação no mesmo servidor, instalação e atualização 4.5 $\rightarrow$ 5.3 e execução nos ambientes de acesso selecionados. Não é possível medir comportamento sob carga, integração com sistemas institucionais reais, compatibilidade com plugins de terceiros, uso ao longo de um semestre nem conformidade legal. Essas limitações influenciam diretamente a priorização.

\subsection{Modelo de qualidade adaptado}

O modelo da ISO/IEC 25010:2023 foi adaptado por quatro regras: (A1)~manter apenas o que apoia uma decisão D1--D3; (A2)~excluir capacidade de interação, por regra da disciplina; (A3)~manter apenas subcaracterísticas mensuráveis na instalação local; e (A4)~reinterpretar cada subcaracterística no vocabulário de um LMS. A Figura~\ref{fig:modelo} mostra o modelo resultante e a Tabela~\ref{tab:subcar}, a interpretação de cada subcaracterística.

\begin{figure}[!ht]
\centering
\begin{tikzpicture}[
  raiz/.style={draw, rounded corners, fill=blue!12, align=center, font=\footnotesize\bfseries, minimum width=4.2cm},
  car/.style={draw, rounded corners, fill=orange!12, align=center, font=\scriptsize\bfseries, minimum width=3.2cm, minimum height=0.8cm},
  sub/.style={draw, rounded corners, fill=white, align=center, font=\scriptsize, minimum width=2.4cm},
  fora/.style={draw, dashed, rounded corners, text=gray, align=center, font=\tiny},
  seta/.style={-{Stealth}, semithick}]
\node[raiz] (q) {Qualidade de produto -- Moodle 5.3 LTS\\(ISO/IEC 25010:2023)};
\node[car, below=0.7cm of q, xshift=-4.3cm] (af) {1. Adequação funcional\\{\tiny risco acadêmico -- D1}};
\node[car, below=0.7cm of q] (co) {2. Compatibilidade\\{\tiny risco de integração -- D2}};
\node[car, below=0.7cm of q, xshift=4.3cm] (fl) {3. Flexibilidade\\{\tiny risco operacional -- D3}};
\foreach \n in {af,co,fl} \draw[seta] (q) -- (\n);
\node[sub, below=0.4cm of af] (af1) {Completude funcional};
\node[sub, below=0.15cm of af1] (af2) {Correção funcional};
\node[sub, below=0.15cm of af2] (af3) {Pertinência funcional};
\node[sub, below=0.4cm of co] (co1) {Interoperabilidade};
\node[sub, below=0.15cm of co1] (co2) {Coexistência};
\node[sub, below=0.4cm of fl] (fl1) {Adaptabilidade};
\node[sub, below=0.15cm of fl1] (fl2) {Instalabilidade};
\node[fora, below=0.15cm of fl2] (fl3) {Escalabilidade e\\substituibilidade (fora)};
\node[fora, below=0.15cm of co2, text width=3.2cm] (out) {Fora: capacidade de interação (A2); eficiência, confiabilidade, segurança, manutenibilidade e segurança operacional (priorização)};
\foreach \a/\b in {af/af1,af1/af2,af2/af3,co/co1,co1/co2,fl/fl1,fl1/fl2} \draw[semithick] (\a) -- (\b);
\draw[semithick, dashed] (fl2) -- (fl3);
\end{tikzpicture}
\caption{Modelo de qualidade adaptado, com a prioridade, o risco e a decisão de cada característica. Fonte: elaboração própria a partir de~\cite{iso25010}.}
\label{fig:modelo}
\end{figure}

\begin{table}[!ht]
\centering
\caption{Subcaracterísticas mantidas e sua interpretação no Moodle.}
\label{tab:subcar}
\footnotesize
\begin{tabularx}{\linewidth}{p{3.1cm}L}
\toprule
\textbf{Subcaracterística} & \textbf{Interpretação no Moodle} \\
\midrule
Completude funcional & O núcleo oferece funções para todas as tarefas do catálogo de tarefas essenciais \\
Correção funcional & Notas, agregações e prazos produzem o resultado esperado pelas regras configuradas \\
Pertinência funcional & As funções levam ao objetivo da tarefa sem exigir contornos \\
Interoperabilidade & Troca de dados por arquivos e serviços REST, com dados utilizáveis pelo receptor \\
Coexistência & Operação no mesmo servidor que outra aplicação e outro banco, sem prejudicá-los \\
Adaptabilidade & Fluxos essenciais funcionam nos navegadores, no celular e nas versões de PHP previstas \\
Instalabilidade & Instalação limpa e atualização 4.5 $\rightarrow$ 5.3 concluem, em tempo razoável e preservando dados \\
\bottomrule
\end{tabularx}
\tabfonte{elaboração própria a partir de~\cite{iso25010,iso25023}.}
\end{table}

\subsection{Seleção e priorização das características}

As oito características restantes (todas, exceto capacidade de interação) receberam notas de 1 a 5 em quatro critérios, com âncoras definidas para cada nota: \textbf{R}~relevância para as partes interessadas, \textbf{K}~risco/impacto de uma falha, \textbf{V}~viabilidade de medição e \textbf{A}~alinhamento às decisões. Foram consideradas elegíveis as características com $V \geq 3$ e $R \geq 2$. A pontuação do cenário base é $P = 0{,}35R + 0{,}25K + 0{,}25V + 0{,}15A$; dois cenários alternativos (viabilidade maior: 0,25/0,20/0,40/0,15; relevância maior: 0,50/0,20/0,15/0,15) testam a sensibilidade. A Tabela~\ref{tab:prior} mostra o resultado.

\begin{table}[!ht]
\centering
\caption{Matriz de priorização (notas, pontuações nos três cenários e elegibilidade).}
\label{tab:prior}
\footnotesize
\begin{tabular}{lcccccccl}
\toprule
\textbf{Característica} & \textbf{R} & \textbf{K} & \textbf{V} & \textbf{A} & \textbf{Base} & \textbf{Viab.} & \textbf{Relev.} & \textbf{Elegível} \\
\midrule
\textbf{Adequação funcional} & 5 & 5 & 4 & 5 & \textbf{4,75} & 4,60 & 4,85 & Sim \\
\textbf{Compatibilidade} & 4 & 4 & 4 & 5 & \textbf{4,15} & 4,15 & 4,15 & Sim \\
\textbf{Flexibilidade} & 4 & 4 & 4 & 4 & \textbf{4,00} & 4,00 & 4,00 & Sim \\
Segurança & 5 & 5 & 2 & 3 & 3,95 & 3,50 & 4,25 & Não ($V<3$) \\
Eficiência de desempenho & 4 & 4 & 2 & 3 & 3,35 & 3,05 & 3,55 & Não ($V<3$) \\
Confiabilidade & 4 & 4 & 2 & 3 & 3,35 & 3,05 & 3,55 & Não ($V<3$) \\
Manutenibilidade & 3 & 3 & 4 & 2 & 3,10 & 3,25 & 3,00 & Sim \\
Segurança operacional & 1 & 1 & 2 & 1 & 1,25 & 1,40 & 1,15 & Não \\
\bottomrule
\end{tabular}
\tabfonte{elaboração própria. As pontuações foram conferidas por script.}
\end{table}

A ordem das elegíveis (adequação funcional $>$ compatibilidade $>$ flexibilidade $>$ manutenibilidade) é a mesma nos três cenários; foram selecionadas as três primeiras. A flexibilidade recebeu $K=4$ porque a mudança da raiz web para \texttt{public/} e os novos requisitos de PHP e PostgreSQL aumentam o risco de falha na atualização. Os principais \emph{trade-offs} são: (i)~a \emph{segurança} teria a segunda maior pontuação sem o filtro de viabilidade, mas exigiria testes de intrusão fora do alcance da equipe, o que é registrado como limitação; (ii)~\emph{eficiência} e \emph{confiabilidade} exigiriam ambiente de carga e observação prolongada; e (iii)~a \emph{manutenibilidade} é viável por análise estática, mas apoia pouco decisões de quem opera, e não de quem desenvolve, o Moodle.

\subsection{Escopo, profundidade e limites de validade}

Os objetos de avaliação são o núcleo do Moodle 5.3 LTS; um catálogo de 12 tarefas essenciais (criar curso e matricular, publicar material, criar e enviar tarefa, tentar envio após o prazo, corrigir, criar e responder questionário, usar fórum, configurar e consultar o livro de notas, fazer backup e restaurar); 7 formatos de arquivo e 6 funções REST; um servidor compartilhado com uma aplicação vizinha; 4 ambientes de acesso (Windows com Chrome e com Firefox, Ubuntu com Firefox, Android com Chrome) e 2 versões de PHP; e os procedimentos oficiais de instalação e atualização. A avaliação é externa, predominantemente caixa-preta e orientada a tarefas. Ficam fora a capacidade de interação, as características não priorizadas, escalabilidade e substituibilidade, plugins de terceiros, o aplicativo móvel, integrações LTI/SCORM/SSO e a instalação institucional em produção. Os resultados valerão apenas para a versão 5.3.x registrada, a configuração de referência e os ambientes testados.

\subsection{Objetivos de Desenvolvimento Sustentável}

A avaliação relaciona-se ao ODS~4, Educação de qualidade~\cite{onu_ods4}. A meta~4.3 (acesso igualitário à educação superior de qualidade; indicador 4.3.1) relaciona-se à cobertura e à correção das tarefas essenciais (M1.1 e M1.2); a meta~4.4 (habilidades relevantes, inclusive em TIC; indicador 4.4.1), de forma indireta, à troca de conteúdos entre cursos (M2.1); e a meta~4.a (ambientes de aprendizagem eficazes; indicador 4.a.1, acesso a internet e computadores para fins pedagógicos), à adaptabilidade a ambientes de acesso e à taxa de sucesso de instalação (M3.1 e M3.3), que reduzem barreiras técnicas para instituições com poucos recursos. Os indicadores oficiais não serão calculados; as métricas locais representam contribuições potenciais, e o funcionamento em celulares não comprova acessibilidade, que pertence à capacidade de interação.

% ======================================================================
\section{Fase 2: especificação da avaliação}
\label{sec:fase2}

\subsection{Objetivos, questões e hipóteses}

Cada característica priorizada gerou um objetivo de medição. Os três objetivos compartilham o objeto (Moodle 5.3 LTS na instalação de referência), o propósito (\emph{entender}) e o ponto de vista (Coordenação de TI), o que garante a rastreabilidade com a Fase~1. As hipóteses vêm das \emph{abstraction sheets} elaboradas pela equipe, que também registraram fatores de variação, como a versão corretiva 5.3.x, a codificação de caracteres dos arquivos CSV e a experiência do avaliador com Linux. A Tabela~\ref{tab:gqm} apresenta os objetivos, as questões e as hipóteses.

\begin{table}[!ht]
\centering
\caption{Objetivos de medição, questões e hipóteses.}
\label{tab:gqm}
\scriptsize
\begin{tabularx}{\linewidth}{p{0.75cm}Lp{4.3cm}}
\toprule
\multicolumn{3}{p{\dimexpr\linewidth-2\tabcolsep}}{\textbf{G1.} \emph{Analisar} o núcleo do Moodle 5.3 LTS \emph{para o propósito de} entender \emph{com respeito à} adequação funcional das tarefas acadêmicas essenciais \emph{do ponto de vista da} Coordenação de TI \emph{no contexto da} decisão de adoção/atualização (D1), em instalação local com dados fictícios.} \\
\midrule
Q1.1 & Em que medida as funções do núcleo cobrem as tarefas essenciais do catálogo? & H1.1: as 12 tarefas são cobertas ($M1.1 = 1{,}00$) \\
Q1.2 & Os resultados calculados (notas, agregação, prazos) correspondem aos esperados? & H1.2: todos os casos iguais ao oráculo ($M1.2 = 1{,}00$) \\
Q1.3 & As funções permitem concluir as tarefas sem contornos? & H1.3: no máximo 1 contorno ($M1.3 \geq 0{,}92$) \\
\midrule
\multicolumn{3}{p{\dimexpr\linewidth-2\tabcolsep}}{\textbf{G2.} \emph{Analisar} o núcleo do Moodle 5.3 LTS e suas interfaces de troca de dados \emph{para o propósito de} entender \emph{com respeito à} compatibilidade (interoperabilidade e coexistência) \emph{do ponto de vista da} Coordenação de TI \emph{no contexto da} decisão sobre as integrações (D2), com sistemas externos simulados.} \\
\midrule
Q2.1 & Em que medida os formatos de arquivo são exportados e importados com sucesso? & H2.1: ao menos 6 de 7 formatos ($M2.1 \geq 0{,}85$) \\
Q2.2 & Os dados trocados chegam íntegros ao outro lado? & H2.2: todos os registros íntegros ($M2.2 = 1{,}00$) \\
Q2.3 & As funções REST selecionadas atendem corretamente a um cliente externo? & H2.3: 6 de 6 funções ($M2.3 = 1{,}00$) \\
Q2.4 & O Moodle opera no mesmo servidor que outras aplicações sem prejudicá-las? & H2.4: nenhuma entidade afetada ($M2.4 = 1{,}00$) \\
\midrule
\multicolumn{3}{p{\dimexpr\linewidth-2\tabcolsep}}{\textbf{G3.} \emph{Analisar} o Moodle 5.3 LTS e seus procedimentos oficiais de instalação e atualização \emph{para o propósito de} entender \emph{com respeito à} flexibilidade (adaptabilidade e instalabilidade) \emph{do ponto de vista da} Coordenação de TI \emph{no contexto da} definição do procedimento de migração (D3), em máquina virtual de referência.} \\
\midrule
Q3.1 & Os fluxos essenciais funcionam em todos os ambientes de acesso previstos? & H3.1: pior ambiente com $M3.1 \geq 0{,}875$ \\
Q3.2 & O Moodle se comporta da mesma forma nas versões de PHP suportadas? & H3.2: sem diferença entre 8.3 e 8.4 ($M3.2 = 1{,}00$) \\
Q3.3 & Instalação limpa e atualização concluem com êxito e em tempo razoável? & H3.3: ao menos uma atualização falha na 1ª tentativa por causa de \texttt{public/} ($M3.3 \geq 0{,}83$); instalação limpa em até 15~min \\
Q3.4 & A atualização 4.5 $\rightarrow$ 5.3 preserva os dados acadêmicos? & H3.4: todos os itens preservados ($M3.5 = 1{,}00$) \\
\bottomrule
\end{tabularx}
\tabfonte{elaboração própria, com o template de objetivo de medição de~\cite{ramos2026aulas}.}
\end{table}

\subsection{Métricas}

As métricas foram derivadas das medidas da ISO/IEC 25023~\cite{iso25023} e seguem a forma de razão, com escala de 0 a 1 em que valores maiores são melhores, exceto M3.4 (tempo). Todas usam método objetivo, com critério de sucesso definido antes da coleta, para atender às propriedades de simplicidade, objetividade, facilidade de coleta e validade~\cite{ramos2026aulas}. A Tabela~\ref{tab:metricas} resume as métricas; a ficha completa de cada uma (definição, medidas básicas, fonte, procedimento de coleta, escala e responsável) está na GitPage.

\begin{table}[!ht]
\centering
\caption{Métricas, fórmulas, medida de referência e fonte dos dados.}
\label{tab:metricas}
\scriptsize
\begin{tabularx}{\linewidth}{p{0.65cm}p{2.25cm}p{1.9cm}p{2.6cm}L}
\toprule
\textbf{ID} & \textbf{Métrica} & \textbf{Fórmula} & \textbf{Medida (25023)} & \textbf{Medidas básicas e fonte} \\
\midrule
M1.1 & Cobertura funcional & $1 - T_f/T_t$ & Functional implementation coverage & $T_t=12$ tarefas; $T_f$ sem função ou com falha; formulário de execução \\
M1.2 & Correção funcional & $C_c/C_t$ & Correctness & $C_t=20$ casos (questionário, agregação, prazo); oráculo calculado à mão \\
M1.3 & Pertinência funcional & $1 - T_c/T_t$ & Functional appropriateness & $T_c$ tarefas que exigiram passo fora da documentação \\
M2.1 & Intercâmbio de formatos & $F_s/F_t$ & Data exchangeability & $F_t=7$: CSV usuários, CSV e XLSX notas, GIFT, Moodle XML, iCal, \texttt{.mbz} \\
M2.2 & Integridade dos dados & $R_c/R_t$ & Data exchangeability (registro) & Registros idênticos à referência em todos os campos \\
M2.3 & Conectividade REST & $W_s/W_t$ & Connectivity with external system & $W_t=6$ funções; HTTP 200, sem exceção, dados iguais à interface \\
M2.4 & Coexistência disponível & $E_c/E_t$ & Available co-existence & $E_t=3$: aplicação web, banco e tarefa \texttt{cron} vizinhos \\
M3.1 & Adaptabilidade ao acesso & $\min_i(1 - T_{fa}/T_{ta})$ & Environmental adaptability & 8 tarefas em cada ambiente A1--A4 \\
M3.2 & Adaptabilidade ao PHP & $\min_i(1 - T_{fp}/T_{tp})$ & System software env. adaptability & 12 tarefas em PHP 8.3 e 8.4 \\
M3.3 & Sucesso de instalação & $P_s/P_t$ & Ease of installation & $P_t=6$ (3 instalações, 3 atualizações) \\
M3.4 & Tempo de instalação & mediana$(t_f-t_i)$; $R$ & Installation time efficiency & Log do terminal; $R$ = retentativas \\
M3.5 & Preservação de dados & $I_p/I_t$ & Definida pela equipe & $I_t=40$ itens contados antes e depois \\
\bottomrule
\end{tabularx}
\tabfonte{elaboração própria a partir de~\cite{iso25023}.}
\end{table}

\subsection{Níveis de pontuação e critérios de julgamento}
\label{sec:julgamento}

Cada métrica é convertida em uma nota de 1 a 4: \textbf{4}~Excelente (nenhuma ação), \textbf{3}~Bom (desvio pequeno com contorno conhecido), \textbf{2}~Insuficiente (exige ação antes da migração) e \textbf{1}~Crítico (impede a tarefa ou corrompe dados). A Tabela~\ref{tab:niveis} define as faixas. M1.2, M2.4 e M3.5 não têm nível~3, pois uma nota final errada, um serviço vizinho afetado ou um envio perdido não são desvios pequenos. Os limites de M3.4 são candidatos e serão calibrados em um piloto na Fase~3, antes da coleta.

\begin{table}[!ht]
\centering
\caption{Níveis de pontuação das métricas.}
\label{tab:niveis}
\scriptsize
\begin{tabularx}{\linewidth}{p{2.4cm}p{1.6cm}p{1.8cm}Lp{2.2cm}}
\toprule
\textbf{Métrica} & \textbf{Nota 4} & \textbf{Nota 3} & \textbf{Nota 2} & \textbf{Nota 1} \\
\midrule
M1.1 Cobertura & 1,00 & $\geq 0{,}92$ & $\geq 0{,}83$ & $< 0{,}83$ \\
M1.2 Correção & 1,00 & -- & $\geq 0{,}95$, só prazo, sem alterar nota & $<0{,}95$ ou nota final errada \\
M1.3 Pertinência & 1,00 & $\geq 0{,}92$ & $\geq 0{,}75$ & $< 0{,}75$ \\
M2.1 Formatos & 1,00 & $\geq 0{,}85$ & $\geq 0{,}70$ & $< 0{,}70$ \\
M2.2 Integridade & 1,00 & $\geq 0{,}99$ & $\geq 0{,}95$ & $< 0{,}95$ \\
M2.3 REST & 1,00 & $\geq 0{,}83$ & $\geq 0{,}66$ & $< 0{,}66$ \\
M2.4 Coexistência & 1,00 & -- & 0,66, com recuperação automática & demais casos \\
M3.1 Acesso (pior) & 1,00 & $\geq 0{,}875$ & $\geq 0{,}75$ & $< 0{,}75$ \\
M3.2 PHP (pior) & 1,00 & $\geq 0{,}92$ & $\geq 0{,}83$ & $< 0{,}83$ \\
M3.3 Sucesso & 1,00 & $\geq 0{,}83$ & $\geq 0{,}66$ & $< 0{,}66$ \\
M3.4 Tempo e $R$ & $\leq$15 min, $R=0$ & $\leq$30 min, $R\leq1$ & $\leq$60 min, $R\leq3$ & $>$60 min ou $R>3$ \\
M3.5 Preservação & 1,00 & -- & $\geq 0{,}975$, item não acadêmico & perda de nota, envio ou usuário \\
\bottomrule
\end{tabularx}
\tabfonte{elaboração própria. Os limites são decisões da equipe, não valores normativos.}
\end{table}

A nota de cada questão é a nota de sua métrica; em Q3.3, é o mínimo entre M3.3 e M3.4. O julgamento do objetivo usa o pior caso entre suas questões, porque uma falha grave em uma subcaracterística não é compensada por bons resultados em outra: se todas as questões têm nota $\geq 3$, o objetivo \textbf{atende} e a decisão correspondente pode prosseguir; se nenhuma tem nota~1 e ao menos uma tem nota~2, \textbf{atende com ressalvas} e a decisão prossegue com ações de mitigação; se alguma tem nota~1, \textbf{não atende} e a falha torna-se candidata à ação de melhoria (D4). Na Fase~4, cada valor também será comparado à hipótese da questão, e hipóteses refutadas serão explicadas pelos fatores de variação das \emph{abstraction sheets}.

\subsection{Diagrama GQM}

A Figura~\ref{fig:gqm} mostra a hierarquia completa nos três níveis do GQM. Toda métrica responde a uma questão, toda questão tem hipótese e todo objetivo deriva de uma característica priorizada e de uma decisão da Fase~1.

\begin{figure}[!ht]
\centering
\begin{tikzpicture}[
  g/.style={draw, rounded corners, fill=orange!15, font=\scriptsize\bfseries, align=center, minimum width=2.1cm, minimum height=0.75cm},
  q/.style={draw, rounded corners, fill=blue!6, font=\tiny, align=center, minimum width=2.4cm, minimum height=0.48cm},
  m/.style={draw, rounded corners, fill=green!7, font=\tiny, align=center, minimum width=2.4cm, minimum height=0.48cm},
  lv/.style={font=\scriptsize\itshape, text=gray, anchor=east}]
\matrix[row sep=0.06cm, column sep=0.8cm] (mx) {
 & \node[q] (q11) {Q1.1 Cobertura}; & \node[m] (m11) {M1.1 Cobertura funcional}; \\
 \node[g] (g1) {G1\\Adequação\\funcional}; & \node[q] (q12) {Q1.2 Correção}; & \node[m] (m12) {M1.2 Correção funcional}; \\
 & \node[q] (q13) {Q1.3 Sem contorno}; & \node[m] (m13) {M1.3 Pertinência}; \\[0.15cm]
 & \node[q] (q21) {Q2.1 Formatos}; & \node[m] (m21) {M2.1 Intercâmbio}; \\
 \node[g] (g2) {G2\\Compatibilidade}; & \node[q] (q22) {Q2.2 Integridade}; & \node[m] (m22) {M2.2 Integridade}; \\
 & \node[q] (q23) {Q2.3 Serviços web}; & \node[m] (m23) {M2.3 Conectividade REST}; \\
 & \node[q] (q24) {Q2.4 Coexistência}; & \node[m] (m24) {M2.4 Coexistência}; \\[0.15cm]
 & \node[q] (q31) {Q3.1 Ambientes de acesso}; & \node[m] (m31) {M3.1 Adaptab. acesso}; \\
 \node[g] (g3) {G3\\Flexibilidade}; & \node[q] (q32) {Q3.2 Versões de PHP}; & \node[m] (m32) {M3.2 Adaptab. PHP}; \\
 & \node[q] (q33) {Q3.3 Instalação}; & \node[m] (m33) {M3.3 Sucesso -- M3.4 Tempo}; \\
 & \node[q] (q34) {Q3.4 Preservação}; & \node[m] (m35) {M3.5 Preservação}; \\
};
\foreach \g/\qs in {g1/{q11,q12,q13}, g2/{q21,q22,q23,q24}, g3/{q31,q32,q33,q34}}
  \foreach \q in \qs \draw[-{Stealth}] (\g.east) -- (\q.west);
\foreach \q/\m in {q11/m11,q12/m12,q13/m13,q21/m21,q22/m22,q23/m23,q24/m24,q31/m31,q32/m32,q33/m33,q34/m35}
  \draw[-{Stealth}] (\q) -- (\m);
\node[font=\scriptsize\bfseries, above=0.15cm of g1 |- q11.north] {Conceitual};
\node[font=\scriptsize\bfseries, above=0.15cm of q11] {Operacional};
\node[font=\scriptsize\bfseries, above=0.15cm of m11] {Quantitativo};
\end{tikzpicture}
\caption{Diagrama GQM com os níveis conceitual (objetivos), operacional (questões) e quantitativo (métricas). Fonte: elaboração própria.}
\label{fig:gqm}
\end{figure}

% ======================================================================
\section{Considerações parciais}
\label{sec:consideracoes}

Esta entrega estabeleceu os requisitos da avaliação do Moodle 5.3 LTS e especificou como ela será feita. A Fase~1 definiu o propósito (apoiar a decisão de migração), os requisitantes e as partes interessadas, o tipo de produto, o modelo de qualidade adaptado, a priorização por método explícito, o escopo e a relação com o ODS~4. A Fase~2 desdobrou as três características priorizadas em 3 objetivos, 11 questões com hipóteses e 12 métricas com níveis de pontuação e critérios de julgamento ligados às decisões do requisitante.

As principais limitações são: as necessidades de várias partes interessadas são hipóteses não validadas por entrevistas; a instalação local não representa uma instalação institucional; a exclusão da segurança depende do filtro de viabilidade; e os limites de tempo de M3.4 ainda precisam ser calibrados. Na Entrega~2, a equipe produzirá o plano de avaliação (Fase~3), com ambiente, massa de dados, oráculo e cronograma; executará as medições e o julgamento (Fase~4); e implementará ao menos uma ação de melhoria justificada pelas evidências.

\bibliographystyle{sbc}
\bibliography{referencias}

\clearpage
\appendix

\section{Tabela de contribuição da equipe}
\label{ap:contribuicao}

\begin{table}[!ht]
\centering
\caption{Contribuição dos integrantes na Entrega 1.}
\label{tab:contribuicao}
\footnotesize
\begin{tabularx}{\linewidth}{p{3.6cm}Lp{1.2cm}p{1.6cm}}
\toprule
\textbf{Integrante} & \textbf{Tarefas realizadas} & \textbf{Horas} & \textbf{\% do esforço} \\
\midrule
Bernardo Broetto Brun & [a preencher pela equipe] & & \\
Carlos Eduardo Deusdara Motta & [a preencher pela equipe] & & \\
Maykon Junio dos Santos Soares & [a preencher pela equipe] & & \\
Pedro Goncalves Rocha & [a preencher pela equipe] & & \\
\midrule
\textbf{Total} & & & \textbf{100\%} \\
\bottomrule
\end{tabularx}
\tabfonte{registro de horas da equipe (\texttt{dados/gestao/horas.csv}).}
\end{table}

\section{Declaração de uso de inteligência artificial}
\label{ap:ia}

[A preencher pela equipe: para cada uso de ferramenta de IA, informar a ferramenta, a tarefa em que foi usada, quando, e como a equipe conferiu e ajustou o resultado. O registro detalhado fica na página ``Declaração de uso de IA'' da GitPage.]

\end{document}
